\documentclass[11pt]{article}
\usepackage[T1]{fontenc}
\usepackage{lmodern}
\usepackage{amsmath}
\usepackage{graphicx}
\usepackage{caption}
\usepackage{hyperref}
\begin{document}

\section{Numbering by hand}
Lists sometimes continue a numbering that was interrupted by a paragraph of prose. The counter commands below are ordinary body text on the fast path: the server restores the counters after every compile and reports what the unit advanced, so the layout pass that renumbers what follows is only scheduled when that changes.

\begin{enumerate}
  \item The first item of the list.
  \item The second item of the list.
\end{enumerate}

An interrupting paragraph of prose between the two parts of the list, long enough to wrap onto a second line so that its own line breaking is exercised as well.

\begin{enumerate}
  \setcounter{enumi}{2}
  \item The third item, numbered by \verb|\setcounter|.
  \item The fourth item follows naturally.
  \addtocounter{enumi}{5}
  \item The tenth item, after \verb|\addtocounter|.
\end{enumerate}

\section{\texorpdfstring{Headings with math $x^2$}{Headings with math}}
The heading above carries a \verb|\texorpdfstring| so the bookmark gets a plain string while the typeset heading keeps the math. The paragraph after it continues as usual.

\stepcounter{equation}
Equations can be skipped with \verb|\stepcounter|; the next numbered equation accounts for it:
\begin{equation}
  e^{i\pi} + 1 = 0 \label{eq:euler}
\end{equation}
and Equation~\eqref{eq:euler} refers to it.

\begin{center}
  \includegraphics[width=0.3\textwidth]{figure.png}
  \captionof{figure}{A figure caption outside a float, with the caption package.}\label{fig:inline}
\end{center}

Figure~\ref{fig:inline} is not a float: its caption is set where it appears, and the figure counter it advances is reported like any other. A number can also be taken inside a paragraph\refstepcounter{equation}\label{eq:taken}: this sentence owns Equation~\eqref{eq:taken}.

\end{document}
